\section{Выводы}

Прямой ход метода Гаусса распараллелен на потоках POSIX: число одновременно
работающих потоков задаётся ключом \texttt{-t} и подтверждается средствами ОС
(\texttt{ps -L}, \texttt{/proc/<pid>/status}, \texttt{top -H}, \texttt{strace}).
Строки распределены между потоками циклически, шаги разделены барьером
\texttt{pthread\_barrier\_t}, общие данные защищены порядком <<запись до
барьера --- чтение после>>. Ответ не зависит от числа потоков. В журнале
автотестов успешно выполнены 46 проверок, включая отказ \texttt{pthread\_create}
без зависания.

Многопоточность окупается только на больших системах: ускорение 1,9 при
$n = 2000$ и 2,2 при $n = 3000$ на 4 потоках. Для $n \le 1000$ засыпание и
пробуждение потоков на $2n$ барьерах дороже полезной работы шага, и программа
в нескольких потоках работает медленнее. Сверху ускорение ограничивают
пропускная способность памяти и число физических ядер.

Ограничения: обратный ход последователен; замеры выполнены в виртуальной машине
WSL и имеют заметный разброс.
